%! Author = Evert
%! Date = 29/09/2023

% Preamble
\documentclass[11pt]{article}

% Packages
\usepackage{amsmath}
\usepackage[top=2cm, bottom=4.5cm, left=2.5cm, right=2.5cm]{geometry}
\usepackage[utf8]{inputenc}
\usepackage[backend=biber]{biblatex}

\addbibresource{references.bib}
\title{Lifelong Learning for an Origami Robot}
\author{Evert De Vroey}

% Document
\begin{document}
    \maketitle


    \section{Introduction}\label{sec:intro}

    \subsection{Challenges}\label{subsec:challenges}
    % catastrophic forgetting
    % handling memories
    % detecting change
    \cite{Kim2021,Lesort2020}.

    \subsection{Application to soft (origami) robots}\label{subsec:origami}
    % reason: change in dynamics
    % no changes in tasks


    \section{Background and definition}\label{sec:background}
    ~\cite{Lesort2020}.


    \section{Blackbox methods}\label{sec:blackbox}
    ~\cite{Zhou2022, Romeres2016, Gijsberts2011, Thuruthel2019, Nguyen2009}.


    \section{Dual architectures}\label{sec:dual}
    ~\cite{Mallya2018, Schwarz2018, Shin2017, Lu2019, Lu2020}.


    \section{Storage and retrieval}\label{sec:storage}
    ~\cite{Wang2022, Schulman2017, Mouret2015, Shin2017, Aljundi2016, Nguyen2011, Xie2021}.


    \section{Architectural expansions}\label{sec:architectural}
    ~\cite{Chatzilygeroudis2016, Nazeer2023, Gillespie2018}.


    \section{Discussion and conclusion}\label{sec:discussion}
    % promising approach(es)
    % research gap

    \printbibliography

\end{document}